% 数学ノート集・共通テンプレート v2（積分論2026年度第1部の版面を参考）
% A5判。章内で定義・定理・例・注意等の番号を共通化。
% コンパイル： platex -interaction=nonstopmode math_note_template_v2.tex (2回)
%              dvipdfmx math_note_template_v2.dvi
\documentclass[a5paper,10pt,onecolumn]{jsarticle}
\usepackage[dvipdfmx,top=22mm,bottom=20mm,inner=19mm,outer=18mm,headheight=13pt,headsep=6mm]{geometry}
\usepackage{amsmath,amssymb,amsthm,mathtools}
\usepackage{enumitem,fancyhdr,titlesec}
\usepackage[dvipdfmx,bookmarks=true,colorlinks=false]{hyperref}
\usepackage{pxjahyper}

% ----- 差し替える文書情報 -----
\newcommand{\NoteTitle}{弱位相と弱*位相}
\newcommand{\NoteSeries}{数学ノート集}

% ----- 文字・組版 -----
% 参考PDFの質素な明朝＋Computer Modernの印象を引き継ぎながら，
% A5判で窮屈にならないよう本文の行間を少し確保する．
\linespread{1.13}
\setlength{\parindent}{1zw}
\setlength{\parskip}{0pt}
\setlength{\abovedisplayskip}{6pt plus 2pt minus 2pt}
\setlength{\belowdisplayskip}{6pt plus 2pt minus 2pt}
\setlength{\abovedisplayshortskip}{3pt plus 2pt minus 1pt}
\setlength{\belowdisplayshortskip}{4pt plus 2pt minus 1pt}
\setlist[enumerate]{topsep=3pt,itemsep=1pt,parsep=0pt,leftmargin=2.4zw}
\titleformat{\section}{\normalfont\normalsize\bfseries}{\thesection}{0.8zw}{}
\titlespacing*{\section}{0pt}{2.0ex plus .4ex}{.9ex}
\titleformat{\subsection}{\normalfont\normalsize\bfseries}{\thesubsection}{0.7zw}{}
\titlespacing*{\subsection}{0pt}{1.5ex plus .3ex}{.6ex}
\raggedbottom

% ----- 柱とページ番号：ページ上部、罫線なし -----
\fancypagestyle{mathnote}{%
  \fancyhf{}%
  \fancyhead[L]{\footnotesize \NoteSeries\ （\NoteTitle）}%
  \fancyhead[R]{\footnotesize \thepage}%
  \renewcommand{\headrulewidth}{0pt}%
  \renewcommand{\footrulewidth}{0pt}%
}
\pagestyle{mathnote}

% ----- 節内通し番号の定義・定理群 -----
\newtheoremstyle{mathnotetheorem}
 {4.5pt}{4.5pt}{\normalfont}{}%
 {\bfseries}{}{0.65zw}%
 {\thmname{#1}\thmnumber{ #2}\thmnote{（#3）}}
\theoremstyle{mathnotetheorem}
\newtheorem{theorem}{定理}[section]
\newtheorem{definition}[theorem]{定義}
\newtheorem{lemma}[theorem]{補題}
\newtheorem{proposition}[theorem]{命題}
\newtheorem{corollary}[theorem]{系}
\newtheorem{example}[theorem]{例}
\newtheorem{remark}[theorem]{注意}
\renewcommand{\proofname}{証明}
\renewcommand{\qedsymbol}{$\square$}
\makeatletter
\renewenvironment{proof}[1][\proofname]{%
  \par\pushQED{\qed}\normalfont\topsep4\p@\@plus2\p@\relax
  \trivlist\item[\hskip\labelsep\bfseries #1]\ignorespaces
}{\popQED\endtrivlist\@endpefalse}
\makeatother

% ----- 使うノートで適宜拡張する記号 -----
\newcommand{\R}{\mathbb{R}}
\newcommand{\norm}[1]{\lVert #1\rVert}
\newcommand{\abs}[1]{\lvert #1\rvert}

\begin{document}
\section*{読了後にできること}
本稿を読むための前提は，ノルム空間と連続双対，Hahn--Banachの定理，初期位相と積位相，コンパクト性である．
Banach--Alaogluの証明ではTychonoffの定理を用い，弱位相と弱*位相を本稿で定義する．
本稿の具体的な到達課題は，弱収束とノルム収束が異なる具体例を証明する．

\section{初期位相と弱位相}

本稿では $X$ を実Banach空間，$X^*$ を連続双対空間とする．
複素Banach空間の場合も，実部を取ることで同様に扱える．

\begin{definition}
集合 $S$ と位相空間族 $Y_i$，
写像 $f_i:S\to Y_i$（$i\in I$）が与えられたとき，
全ての $f_i$ を連続にする最弱の位相を
\emph{初期位相}という．
この位相の開集合の基は，
有限個の $i_1,\dots,i_n$ と
$Y_{i_j}$ の開集合 $U_j$ に対し
$\bigcap_{j=1}^nf_{i_j}^{-1}(U_j)$ の形で与えられる．
\end{definition}

\begin{definition}
$X$ 上で全ての $f\in X^*$ を連続にする最弱の位相を
\emph{弱位相}といい，$\sigma(X,X^*)$ と書く．
点 $x_0\in X$ における基本近傍は
\[
 U(x_0;f_1,\dots,f_n;\varepsilon)
 =\{x\in X:|f_j(x-x_0)|<\varepsilon\
  (j=1,\dots,n)\}
\]
の形で与えられる．
この位相での収束を弱収束といい，
$x_\alpha\rightharpoonup x$ と書く．
\end{definition}

\begin{proposition}[弱収束の特徴づけ]\label{prop:weak-conv}
$X$ のネット $(x_\alpha)$ が $x$ に弱収束することと
\[
 f(x_\alpha)\to f(x)\qquad(f\in X^*)
\]
が成り立つことは同値である．
特にノルム収束すれば弱収束する．
\end{proposition}
\begin{proof}
弱位相は $f\in X^*$ による初期位相だから，
全ての写像 $f$ を通じた像が収束することが
初期位相での収束の必要十分条件である．
ノルム収束の場合は
$|f(x_\alpha-x)|\le\norm f\,\norm{x_\alpha-x}\to0$
なので後半が従う．
\end{proof}

\begin{proposition}[弱位相のHausdorff性]\label{prop:weak-hausdorff}
弱位相 $\sigma(X,X^*)$ はHausdorffである．
また $X$ が無限次元なら弱位相はノルム位相より真に弱い．
\end{proposition}
\begin{proof}
$x\ne y$ ならHahn--Banachの定理によって
$f\in X^*$ で $f(x-y)\ne0$ となるものが存在する．
$f(x)$ と $f(y)$ の互いに交わらない開近傍を
$\mathbb R$ に取って逆像を考えれば，
$x$ と $y$ は弱開集合で分離される．

$X$ が無限次元とする．任意の弱基本近傍
$U(0;f_1,\dots,f_n;\varepsilon)$ に対して
\[
 M=\bigcap_{j=1}^n\ker f_j
\]
は有限余次元部分空間なので $M\ne\{0\}$ である．
従って非零の $v\in M$ を取ると
全ての $tv$（$t\in\mathbb R$）が
この近傍に属する．
ゆえに弱基本近傍はいずれもノルムについて非有界であり，
ノルム単位球に含まれる弱近傍は存在しない．
従って弱位相とノルム位相は一致しない．
\end{proof}

\section{弱*位相}

\begin{definition}
双対空間 $X^*$ 上で
\[
 \operatorname{ev}_x:X^*\to\mathbb R,\qquad
 \operatorname{ev}_x(f)=f(x)
 \quad(x\in X)
\]
を全て連続にする最弱の位相を
\emph{弱*位相}といい，$\sigma(X^*,X)$ と書く．
\end{definition}

\begin{proposition}\label{prop:weak-star}
$X^*$ のネット $f_\alpha$ が $f$ に弱*収束することと
\[
 f_\alpha(x)\to f(x)\qquad(x\in X)
\]
が成り立つことは同値である．
また弱*位相はHausdorffであり，
$X^*$ のノルム位相より弱い．
\end{proposition}
\begin{proof}
収束の特徴づけは初期位相の定義による．
$f\ne g$ ならある $x\in X$ で $f(x)\ne g(x)$ なので，
評価写像の逆像によって二点を分離できる．
また
$|f_\alpha(x)-f(x)|\le\norm{f_\alpha-f}\norm x$
なのでノルム収束は弱*収束を導く．
\end{proof}

\begin{remark}
$X^*$ には，弱位相 $\sigma(X^*,X^{**})$ と
弱*位相 $\sigma(X^*,X)$ の二つが存在する．
後者では $X$ を自然に $X^{**}$ へ埋め込んで
評価汎関数を指定している．
一般に両位相は一致しない．
\end{remark}

\section{凸集合と弱閉集合}

\begin{theorem}[閉凸集合の弱閉性]\label{thm:weak-closed}
$C\subset X$ を凸集合とする．
このとき $C$ がノルム閉であることと，
弱位相 $\sigma(X,X^*)$ で閉であることは同値である．
\end{theorem}
\begin{proof}
弱位相はノルム位相より弱いので，
弱閉集合はノルム閉である．
逆に $C$ をノルム閉凸集合とする．
空集合の場合は自明である．$x_0\notin C$ を取ると
閉凸集合と外点の強分離定理によって
$f\in X^*$ と $\alpha\in\mathbb R$ が存在して
\[
 f(c)\le\alpha<f(x_0)\qquad(c\in C)
\]
となる．
従って
$V=\{x\in X:f(x)>\alpha\}$ は
$x_0$ を含む弱開集合であり，$C$ と交わらない．
$x_0\notin C$ は任意なので $X\setminus C$ は弱開，
すなわち $C$ は弱閉である．
\end{proof}

\begin{corollary}\label{cor:weak-closure}
任意の凸集合 $C\subset X$ について，
その弱閉包とノルム閉包は一致する．
\end{corollary}
\begin{proof}
ノルム閉包 $\overline C^{\norm{\cdot}}$ は凸かつノルム閉なので，
定理\ref{thm:weak-closed}から弱閉である．
従って弱閉包はノルム閉包に含まれる．
逆に弱位相はノルム位相より弱いので，
ノルム閉包は弱閉包に含まれる．
よって一致する．
\end{proof}

\section{Banach--Alaogluの定理}

\begin{theorem}[Banach--Alaoglu]\label{thm:alaoglu}
$X$ を実ノルム空間とする．
双対空間の閉単位球
\[
 B_{X^*}=\{f\in X^*:\norm f\le1\}
\]
は弱*位相 $\sigma(X^*,X)$ に関してコンパクトである．
\end{theorem}
\begin{proof}
各 $x\in X$ に対して閉区間
$I_x=[-\norm x,\norm x]$ を取る．
積空間
\[
 K=\prod_{x\in X}I_x
\]
は，各 $I_x$ がコンパクトであるから，
Tychonoffの積定理によって積位相に関してコンパクトである．
写像
\[
 J:B_{X^*}\to K,\qquad J(f)=(f(x))_{x\in X}
\]
を考える．$J$ は単射である．
また $B_{X^*}$ の弱*位相は評価写像の初期位相なので，
$J$ はその像への同相写像である．

$J(B_{X^*})$ が $K$ の閉部分集合であることを示す．
$K$ の点 $u=(u_x)_{x\in X}$ が $J(B_{X^*})$ に属することは，
\[
 u_{x+y}=u_x+u_y,\qquad
 u_{ax}=a\,u_x
 \quad(x,y\in X,\ a\in\mathbb R)
 \tag{$*$}
\]
が成立することと同値である．
実際 $(*)$ が成り立つなら $f_u(x)=u_x$ は線形であり，
$u_x\in I_x$ だから
$|f_u(x)|\le\norm x$ を満たして $\norm{f_u}\le1$ である．
逆に $\norm f\le1$ の線形汎関数は $(*)$ を満たす．
$(*)$ は座標射影の有限個の連続演算による
閉条件であり，その全体は閉集合の共通部分なので
$J(B_{X^*})$ は閉である．
コンパクト空間 $K$ の閉部分集合はコンパクトなので，
$J(B_{X^*})$ はコンパクトである．
同相写像 $J$ によって $B_{X^*}$ も弱*コンパクトとなる．
\end{proof}

\begin{remark}
この証明ではTychonoffの積定理を用いた．
これは任意個のコンパクト空間の積がコンパクトであるという
一般位相空間論の定理であり，選択公理と関係する．
Banach--Alaogluの定理の証明が
有限次元のHeine--Borelの定理だけで済むわけではない．
\end{remark}

\section{可分空間の場合の距離化}

\begin{theorem}[弱*単位球の距離化]\label{thm:metrizable}
$X$ を可分ノルム空間とする．
このとき $B_{X^*}$ の弱*位相は距離化可能である．
従って任意の有界列 $(f_n)\subset X^*$ は，
弱*収束する部分列を持つ．
\end{theorem}
\begin{proof}
$X$ の稠密な可算集合を
$\{x_1,x_2,\dots\}$ とする．
$B_{X^*}$ 上で
\[
 d(f,g)=\sum_{j=1}^\infty2^{-j}
 \frac{|f(x_j)-g(x_j)|}{1+|f(x_j)-g(x_j)|}
\]
と定める．各項は非負で1以下なので級数は収束する．
また $d(f,g)=0$ なら
$f(x_j)=g(x_j)$ が全ての $j$ で成り立ち，
連続性と稠密性から $f=g$ である．
$\rho(a,b)=|a-b|/(1+|a-b|)$ は
三角不等式を満たす（実数に対する
$t\mapsto t/(1+t)$ の劣加法性を用いる）ので，
$d$ は距離である．

弱*収束 $f_\alpha\to f$ なら
各 $x_j$ で収束するから，級数の尾部を先に小さく取って
有限個の初項で収束を確認することで
$d(f_\alpha,f)\to0$ を得る．
逆に $d(f_\alpha,f)\to0$ なら
全ての $j$ に対して
$f_\alpha(x_j)\to f(x_j)$ となる．
任意の $x\in X$ と $\varepsilon>0$ に対し
$\norm{x-x_j}<\varepsilon$ となる $x_j$ を取ると，
単位球上で
\[
 |f_\alpha(x)-f(x)|
 \le2\varepsilon+|f_\alpha(x_j)-f(x_j)|
\]
となる．$\varepsilon\downarrow0$ により
全ての $x$ について収束し，弱*収束を得る．

Banach--Alaogluの定理から $B_{X^*}$ は弱*コンパクトであり，
距離化可能なコンパクト空間は点列コンパクトなので
単位球内の任意の列は弱*収束部分列を持つ．
一般の有界列はある $R B_{X^*}$ に含まれるので，
スカラー倍の同相性を使えば結論が従う．
\end{proof}

\section{Goldstineの定理と反射性}

\begin{definition}
自然な写像
\[
 \iota:X\to X^{**},\qquad \iota(x)(f)=f(x)
\]
を\emph{標準埋め込み}という．
$\iota$ が全射であるとき $X$ を\emph{反射的}という．
\end{definition}

\begin{lemma}\label{lem:canonical-isometry}
$\iota$ は線形等距離埋め込みである．
\end{lemma}
\begin{proof}
線形性は定義から従う．
$\norm{\iota(x)}
 =\sup_{\norm f\le1}|f(x)|\le\norm x$ である．
$x\ne0$ ならHahn--Banachの定理から
$\norm f=1$ かつ $f(x)=\norm x$ を満たす
$f\in X^*$ が存在する．
これを代入すると逆向きの不等式が従い，
$\norm{\iota(x)}=\norm x$ となる．
$x=0$ の場合も明らかである．
\end{proof}

\begin{theorem}[Goldstineの定理]\label{thm:goldstine}
$\iota(B_X)$ は $B_{X^{**}}$ において
弱*稠密である．ここで $X^{**}$ には
$\sigma(X^{**},X^*)$ を入れる．
\end{theorem}
\begin{proof}
$\Phi\in B_{X^{**}}$ を取り，
有限個の $f_1,\dots,f_n\in X^*$ と $\varepsilon>0$ を固定する．
$\Phi$ の弱*基本近傍を考えるため，
\[
 T:X\to\mathbb R^n,\qquad
 T(x)=(f_1(x),\dots,f_n(x))
\]
と置く．
$v=(\Phi(f_1),\dots,\Phi(f_n))$ が
$T(B_X)$ の閉包に属することを示す．
もし属さなければ，
有限次元空間の閉凸集合と外点の強分離定理によって
係数 $a_1,\dots,a_n$ が存在して
\[
 \sum_{j=1}^na_j\Phi(f_j)
 >\sup_{\norm x\le1}\sum_{j=1}^na_j f_j(x).
\]
右辺は
$\norm{\sum_ja_jf_j}$ である．
一方，$\norm\Phi\le1$ なので
\[
 \sum_ja_j\Phi(f_j)
 =\Phi\left(\sum_ja_jf_j\right)
 \le\norm{\sum_ja_jf_j},
\]
矛盾する．
従って $v\in\overline{T(B_X)}$ であり，
ある $x\in B_X$ を選んで
$|f_j(x)-\Phi(f_j)|<\varepsilon$（全ての $j$）
とできる．
これは任意の弱*基本近傍が $\iota(B_X)$ と交わることを意味する．
ゆえに弱*稠密である．
\end{proof}

\begin{theorem}[弱コンパクト性と反射性]\label{thm:reflexive}
Banach空間 $X$ について，次は同値である．
\begin{enumerate}
\item $X$ は反射的である．
\item $B_X$ は弱位相 $\sigma(X,X^*)$ でコンパクトである．
\end{enumerate}
\end{theorem}
\begin{proof}
(1) とする．標準埋め込み $\iota$ は
$X$ と $X^{**}$ の間の線形等距離全単射となる．
従って $\iota(B_X)=B_{X^{**}}$ である．
$B_{X^{**}}$ は
$X^*$ の双対の閉単位球なので，
Banach--Alaogluの定理によって
$\sigma(X^{**},X^*)$ に関してコンパクトである．
$\iota$ は弱位相と弱*位相の間の同相写像なので，
$B_X$ は弱コンパクトとなる．

逆に (2) とする．
$\iota:B_X\to B_{X^{**}}$ は
弱位相から弱*位相への連続写像である．
従って $\iota(B_X)$ は
弱*コンパクト空間 $B_{X^{**}}$ のコンパクト部分集合となる．
弱*位相はHausdorffなので，
コンパクト部分集合 $\iota(B_X)$ は閉である．
一方 Goldstineの定理から
$\iota(B_X)$ は $B_{X^{**}}$ に弱*稠密である．
閉かつ稠密なので
$\iota(B_X)=B_{X^{**}}$ となる．
任意の $\Phi\in X^{**}$ をノルムで割って単位球へ移せば
$\Phi$ は $\iota(X)$ に属するので $\iota$ は全射であり，
$X$ は反射的である．
\end{proof}

\section{具体例}

\begin{example}[$\ell^2$ の標準基底の弱収束]
$\ell^2$ の標準基底 $e_n$ は零に弱収束するが，
ノルム収束しない．
\end{example}
\begin{proof}
Rieszの表現定理から
$\ell^2$ 上の任意の連続線形汎関数は
$f_y(x)=\sum_kx_k\overline{y_k}$，
$y\in\ell^2$ の形に表される．
$f_y(e_n)=\overline{y_n}\to0$ は
$\ell^2$ 数列が零に収束することから従う．
従って $e_n\rightharpoonup0$ である．
一方 $\norm{e_n}=1$ なのでノルム収束しない．
\end{proof}

\begin{example}[$\ell^1$ と $\ell^\infty$ の対比]
$\ell^1$ の標準基底 $e_n$ は零に弱収束しない．
一方，$\ell^\infty=(\ell^1)^*$ とみなしたとき，
$\ell^\infty$ の標準基底 $e_n$ は零に弱*収束する．
\end{example}
\begin{proof}
$\ell^1$ 上の汎関数
$f(x)=\sum_kx_k$ は
$|f(x)|\le\norm x_1$ なので連続であり，
$f(e_n)=1$ は零へ収束しない．
従って $\ell^1$ の $e_n$ は零に弱収束しない．

次に $\ell^\infty$ の元としての $e_n$ を考える．
任意の $x=(x_k)\in\ell^1$ に対して
$\langle e_n,x\rangle=x_n\to0$ となる．
これは弱*収束の定義なので，
$\ell^\infty$ の $e_n$ は零に弱*収束する．
\end{proof}

\begin{remark}
弱収束は全ての連続線形汎関数による評価の収束であり，
弱*収束は双対空間の評価写像だけによる収束である．
Banach--Alaogluの定理は弱*コンパクト性を保証するが，
一般のBanach空間の単位球の弱コンパクト性は
反射性と同値であり，別の性質である．
\end{remark}

\section{到達課題の完全解答}
\noindent\textbf{標準直交基底の弱収束．}
\begin{proposition}
Hilbert空間 $\ell^2$ の標準単位ベクトル $e_n$ は
弱く零に収束するが，ノルムでは零に収束しない．
\end{proposition}
\begin{proof}
任意の $y=(y_k)\in\ell^2$ に対し
$\langle e_n,y\rangle=\overline{y_n}$ である．
$\sum_k|y_k|^2<\infty$ だから $y_n\to0$ であり，
$\langle e_n,y\rangle\to0$．
Hilbert空間のRiesz表現定理により
任意の連続線形汎関数は内積で表されるので
$e_n\rightharpoonup0$ である．
一方 $\|e_n-0\|_2=1$ は一定である．
\end{proof}
\noindent\textbf{到達確認．}
$e_n/n$ はノルム収束するか．
\par\noindent\textbf{解答．}
$\|e_n/n\|_2=1/n\to0$ だからノルムで零へ収束する．
\end{document}